\section{Software architecture}
\label{sec:architecture}

Figure~\ref{fig:architecture} shows how \phynetpy\ is put together. The
organising decision is that no analysis owns the network. A network is a data
structure that the library provides; scoring criteria, generative models,
topology moves and search drivers are separate things that act on it. Almost
everything else in this section follows from that separation.

\begin{figure}[t]
  \centering
  \includegraphics[width=\linewidth]{architecture.pdf}
  \caption{\phynetpy\ organised by dependency. A run is specified by picking
  one object from each of the three axes at the top and calling one of the
  three verbs. The registry resolves that triple to an engine. Engines do not
  implement their own search machinery, topology moves, parameter optimisation
  or network operations; they compose the shared components below them. The
  compiled kernels sit under the network layer and are not part of the public
  interface.}
  \label{fig:architecture}
\end{figure}

\subsection{The network as a reusable object}

\code{Network} is a rooted directed acyclic graph over \code{Node} and
\code{Edge}, with \code{Branch} carrying length, inheritance probability and
tag. Trees are not a separate type: a tree is a network with no reticulation
nodes, and the same code paths serve both. Gene trees supplied as input are
\code{Network} instances, which is why a gene tree and an inferred species
network can be passed to the same distance function.

Reticulation in-degree is not capped at two. \code{Network.get_parents}
imposes no limit and \code{add_edges} does not check parent counts, so a node
representing an unresolved or aggregated reticulation with three or more
parents is representable. Some analyses do assume the binary case --
\code{reticulation_tripartitions} requires exactly two parents, for instance --
so the generality lives in the representation rather than throughout the
library.

The operations on \code{Network} are the ones network methods keep needing:
topological ordering, breadth- and depth-first traversal, most recent common
ancestors, leaf-descendant sets for all nodes at once, induced subnetworks,
acyclicity checking, and copying. Ploidy is first class through
\code{subgenome_count} and the corresponding edge queries, and conversion to a
multiply-labelled (MUL) tree is implemented because the allopolyploid parsimony
criterion is defined on that representation.

Two design choices in this layer are worth stating because they change what
callers can get wrong. Branch lengths carry an explicit unit --
\code{SUBSTITUTIONS_PER_SITE}, \code{COALESCENT_2N}, or
\code{UNSPECIFIED} -- and \code{convert_network_branch_lengths} refuses to
convert when either side is unspecified rather than guessing. Coalescent
methods and sequence methods read branch lengths on different scales, and
mixing them produces plausible numbers rather than an error, so the unit is
tracked on the network instead of being left to the user's discipline. Second,
the node and edge containers are compiled (\code{graph_core_cy}) and maintain
indices for incidence and name lookup, because network search evaluates large
numbers of candidate topologies and the cost of a single edge insertion sets
the ceiling on tractable problem size.

\subsection{Three axes and three verbs}

A phylogenetic network method is usually presented as a named procedure. Seen
from the software side, most of those names vary along three independent
dimensions: the data observed, the generative process assumed, and the
statistical objective optimised. \phynetpy\ makes those three dimensions
explicit arguments and reduces the interface to three verbs.

\code{infer} searches for a network. \code{score} evaluates one. \code{simulate}
runs the axes backwards, taking a model and a network and returning a data
object. The axes are:

\begin{itemize}
  \item \textbf{Data} --- \code{GeneTrees}, \code{Alignment},
    \code{BiallelicMarkers}. Each carries the species-to-allele mapping, and
    each knows its own type, so a caller never declares what kind of data it
    is passing.
  \item \textbf{Model} --- \code{MSC}, the multispecies network coalescent with
    a population parameter $\theta = 4N\mu$, and \code{Allopolyploid}, which
    takes a subgenome map. The generative process used to be implicit in which
    scorer a user selected.
  \item \textbf{Criterion} --- \code{MDC} (minimising deep coalescences),
    \code{Likelihood}, \code{PseudoLikelihood}, and \code{Bayesian}.
\end{itemize}

\code{Bayesian} wraps an objective rather than sitting beside the likelihoods:
\code{Bayesian(objective=Likelihood())}. Sampling a posterior is a mode of
working with a likelihood, not a fourth thing to optimise, and expressing it
that way means \code{Bayesian} inherits the data types its objective accepts
instead of restating them.

The practical consequence is that changing method means changing an argument.
The same call that maximises a pseudo-likelihood samples a posterior if the
criterion changes, and \code{infer} returns the same type either way:
\code{InferenceResult} carries \code{.best}, \code{.score}, \code{.trace}, and
\code{.posterior}, the last populated only for Bayesian runs. Attribute access
falls through to \code{.raw}, the engine's native result, so per-method
richness -- log writers, information criteria, reticulation posteriors --
remains reachable without the wrapper having to enumerate it.

\subsection{Dispatch as a validity matrix}

Whether a run makes sense depends on the whole triple, not on any one axis. A
parsimony criterion over a sequence alignment is not a missing feature; it is a
category error, since minimising deep coalescences is defined over gene-tree
topologies. Maximum parsimony from gene trees under the \code{MSC}, by
contrast, is perfectly well defined and simply has no implementation here yet.
A criterion that needs branch lengths applied to gene trees that lack them is a
third, different problem.

\code{resolve} in \code{_registry.py} separates these three cases and raises
accordingly: \code{TypeError} when the criterion is not defined on that data
type, \code{ValueError} when the data lack branch lengths the criterion
requires, and \code{NotImplementedError} when the combination is meaningful but
unbuilt. Conflating them is what makes an unsupported request hard to act on:
a user cannot tell whether to supply different data, wait for a release, or
abandon the idea.

Because dispatch is a registry keyed on the triple, the set of available
methods is introspectable. \code{validity_matrix()} tabulates it, and the table
is built from the registry rather than maintained alongside it, so it cannot
fall out of date. \code{registered_cells()} lists every implemented
combination; Table~\ref{tab:analyses} is generated from that call.

\subsection{Extension}

Adding a method means registering a cell:

\begin{lstlisting}
@register(GeneTrees, MSC, MyCriterion)
class MyEngine(Engine):
    method = "MyMethod"

    def infer(self, data): ...
    def score(self, network, data, optimize=False): ...
\end{lstlisting}

What an engine author does not write is the rest of the run. The search drivers
(\code{HillClimbing}, \code{SimulatedAnnealing}, and the \code{ProposalKernel}
interface), the ten topology-move classes, the Brent coordinate-ascent
optimiser for branch lengths and inheritance probabilities, the I/O layer and
the chain diagnostics are shared. This is the claim in the paper's title that
rests on design rather than measurement, so it is worth being precise about the
evidence: the registry is exercised by all seven of the library's own methods,
which is a weaker statement than demonstrating that an outside contributor
found it sufficient.

One cross-cutting option illustrates the difference between a constraint and a
hint. A starting phylogeny can be supplied as \code{Start(net)}, a free
starting point, or as \code{Start(net, StartMode.AUGMENT)}, which requires the
result to contain that network. The second is enforced by a cluster-containment
check in the accept path rather than approximated by withholding moves, which
is possible because adding reticulations to a network preserves its clusters.
Engines whose proposal kernels cannot honour it -- \code{MCMC_SEQ} and the
allopolyploid search -- reject the request explicitly instead of accepting it
and quietly doing something else.

\subsection{Compiled kernels}

Four Cython modules are compiled at install time. \code{graph_core_cy} provides
the node and edge containers; \code{gt_msc_cy} and \code{mpl_engine_cy}
accelerate the gene-tree coalescent likelihood and the pseudo-likelihood
dynamic program. The fourth, \code{seq_engine_cy}, is compiled but not yet
called from \code{_seq_likelihood.py}.

None of these are part of the public interface, and the graph core is a hard
requirement rather than an optional accelerator: \code{Network} imports it at
module load and raises with build instructions if it is absent. That removes a
class of bug in which a user unknowingly runs a slower parallel implementation,
at the cost of requiring a C compiler when installing from source.
Section~\ref{sec:validation} reports what the compiled pseudo-likelihood kernel
is worth against the pure-Python path that remains in the repository as a
reference.
